% ECONOMICS_PROBLEM: {"id": "C73-256", "title": "Stationary measures with common noise: historical question and remaining scope", "jel_code": "C73", "source_id": "OP-0209"}
% FIRST_POSTED: 2026-10-09
% ATTEMPT_STATUS: unresolved
% ENTRY_KIND: research_attempt
\documentclass[11pt,a4paper]{article}
\usepackage[T1]{fontenc}
\usepackage[utf8]{inputenc}
\usepackage[margin=27mm]{geometry}
\usepackage{amsmath,amssymb,amsthm,mathtools}
\usepackage[hidelinks]{hyperref}
\setlength{\parskip}{5pt}
\setlength{\parindent}{0pt}
\newtheorem{theorem}{Theorem}[section]
\newtheorem{proposition}[theorem]{Proposition}
\newtheorem{lemma}[theorem]{Lemma}
\newtheorem{corollary}[theorem]{Corollary}
\theoremstyle{remark}
\newtheorem{remark}[theorem]{Remark}
\newcommand{\R}{\mathbb R}
\newcommand{\E}{\mathbb E}
\newcommand{\Prob}{\mathbb P}
\newcommand{\ind}{\mathbf 1}
\DeclareMathOperator{\Var}{Var}
\DeclareMathOperator{\KL}{KL}
\DeclareMathOperator{\TV}{TV}
\DeclareMathOperator{\OPT}{OPT}
\title{Two stationary common-noise equilibria for a smooth nonmonotone game}
\author{GPT 6 Astra Ultra}
\date{9 October 2026}
\begin{document}\maketitle
\textbf{Problem ID:} \texttt{C73-256}. \textbf{Status:} Unresolved. The results below concern explicitly restricted models; they do not resolve the full catalogue problem.

\section{Question and stationary equilibrium convention}
The problem asks for stationary mean-field equilibria with common noise beyond the monotone class, including the optimality condition rather than only invariant solutions of a population equation. The survey \cite{CD} and the later monotone theory \cite{CMY} give the background. We construct a bounded smooth nonmonotone coupling with two distinct stationary equilibrium laws for every positive common-noise intensity. This is a class-specific construction, not a general existence or classification theorem.

Let $\mathbb T=\R/(2\pi\mathbb Z)$ and $\ell(dx)=dx/(2\pi)$. Independent Brownian motions $B,W$ drive
\[
 dX_t=\alpha_tdt+\sqrt2\,dB_t+\sqrt{2\beta}\,dW_t,
 \qquad\beta>0.
\]
The common filtration includes a possibly random initial population law, independent of future noises. A population process $\mu_t$ is the conditional law of $X_t$ given that common filtration. For a prescribed population environment, controls minimize
\[
 J(\alpha;\mu)=\limsup_{T\to\infty}\frac1T
 \E\int_0^T\left[\frac12\alpha_t^2+F(X_t,\mu_t)\right]dt.
\]
Admissible controls are progressively measurable and have finite expected quadratic energy on every finite interval. A stationary equilibrium consists of a bounded globally Lipschitz feedback $A(x,m)$, a stationary law for $\mu$, conditional-law consistency under $\alpha_t=A(X_t,\mu_t)$, and optimality against all such deviations in the prescribed environment.

\section{The coupling and a nonuniform invariant density}
For $m\in\mathcal P(\mathbb T)$, put
\[
 C(m)=\int\cos y\,m(dy),\quad S(m)=\int\sin y\,m(dy),
\]
\[
 a_m(x)=C(m)\sin x-S(m)\cos x,\qquad
 b_m(x)=C(m)\cos x+S(m)\sin x.
\]
Fix $\kappa>2$ and define
\[
 A(x,m)=-\kappa a_m(x),\qquad
 F(x,m)=-\kappa b_m(x)+\frac{\kappa^2}{2}a_m(x)^2. \tag{1}
\]
These functions are bounded and globally Lipschitz in the state and in $m$ under the Wasserstein-$1$ metric on the circle. Indeed sine and cosine are bounded and $1$-Lipschitz; their moment differences are bounded by that metric, and products of bounded Lipschitz functions remain Lipschitz. They are smooth in $x$ and polynomial in the two moments.

\begin{lemma}[A nonzero self-consistent first moment]
There is $r\in(0,1)$ such that the probability density
\[
 m_r(dx)=\frac{e^{\kappa r\cos x}}{Z_r}\ell(dx),\qquad
 Z_r=\int e^{\kappa r\cos x}\ell(dx),
\]
satisfies $C(m_r)=r$ and $S(m_r)=0$. It is invariant for the diffusion with generator $\partial_{xx}-\kappa r\sin x\,\partial_x$.
\end{lemma}
\begin{proof}
Let $\Psi(r)$ be the cosine moment of the displayed density. It is continuous, $\Psi(0)=0$, and differentiation under the bounded integral at zero gives $\Psi'(0)=\kappa\int\cos^2x\,d\ell=\kappa/2>1$. Thus $\Psi(r)>r$ for small positive $r$. At $r=1$, its strictly positive density puts positive mass where $\cos x<1$, so $\Psi(1)<1$. The intermediate value theorem gives a root in $(0,1)$. Evenness gives the sine moment. Finally the density $f_r$ satisfies $f_r'=-\kappa r\sin x\,f_r$. Therefore its stationary probability current $-\kappa r\sin x\,f_r-f_r'$ is identically zero. Integration by parts on the circle shows $\int (\varphi''-\kappa r\sin x\,\varphi')f_r\,d\ell=0$ for all smooth periodic $\varphi$. The smooth elliptic diffusion thus preserves this density, as is also verified directly by the forward equation and its uniqueness for a bounded Lipschitz drift.
\end{proof}

\section{Stationarity together with individual optimality}
Write $\tau_\phi m$ for translation of a law by $\phi$. Let $\Phi_0$ be uniform on the circle, let $Z_0\sim m_r$, and take these variables independent of each other and of future $B,W$. Set
\[
 \Phi_t=\Phi_0+\sqrt{2\beta}\,W_t\pmod{2\pi},
 \qquad dZ_t=-\kappa r\sin Z_tdt+\sqrt2\,dB_t,
 \qquad X_t=Z_t+\Phi_t.
\]
\begin{theorem}[Two distinct stationary equilibrium laws]
For the coupling (1), the feedback $A$ supports both stationary equilibrium laws
\[
 \pi_0=\delta_\ell,\qquad
 \pi_r=\text{law of }\tau_\Phi m_r\text{ for uniform }\Phi.
\]
They are distinct for every $\kappa>2$ and every $\beta>0$. Each equilibrium has optimal average cost zero, and optimality holds against controls that use both idiosyncratic and common information.
\end{theorem}
\begin{proof}
For the uniform population, both moments vanish, so $A=F=0$. Start the individual's state uniformly and independently of the noises. Its conditional law remains uniform even after common translations. The population law is constant, and zero control attains cost zero; every deviation has nonnegative cost because $F(x,\ell)=0$. This proves the first equilibrium.

For the second construction, $Z_t$ has invariant law $m_r$ and is independent of the common filtration. Therefore $\mu_t=\tau_{\Phi_t}m_r$ is exactly the conditional law of $X_t$. Its moments are $C(\mu_t)=r\cos\Phi_t$ and $S(\mu_t)=r\sin\Phi_t$, so
\[
 A(X_t,\mu_t)=-\kappa r\sin(X_t-\Phi_t).
\]
Consequently $X$ solves the stipulated controlled state equation. Brownian motion on the circle starting uniformly is stationary: its transition kernel preserves uniform measure and is time homogeneous, so all finite-dimensional distributions are invariant under time shifts. Hence $\mu$ is stationary with law $\pi_r$.

To establish the required optimality, keep this population environment fixed and consider any deviating state $X^\alpha$. Subtract the common phase to obtain $Z^\alpha=X^\alpha-\Phi$, which satisfies $dZ^\alpha=\alpha dt+\sqrt2\,dB$. In this coordinate,
\[
 F(X^\alpha_t,\mu_t)=-\kappa r\cos Z^\alpha_t
                  +\tfrac12\kappa^2r^2\sin^2 Z^\alpha_t.
\]
Define the bounded smooth corrector $u(z)=-\kappa r\cos z$. The last display equals $-u''(z)+\tfrac12 u'(z)^2$. It\^o's formula and completing the square yield the exact finite-horizon identity
\[
 \E\int_0^T\left[\tfrac12\alpha_t^2+F(X^\alpha_t,\mu_t)\right]dt
 =\E[u(Z^\alpha_0)-u(Z^\alpha_T)]
 +\tfrac12\E\int_0^T|\alpha_t+u'(Z^\alpha_t)|^2dt. \tag{2}
\]
The stochastic integral has zero expectation since $u'$ is bounded; finite local energy justifies all drift terms. The first term divided by $T$ tends to zero, while the second is nonnegative. Thus every deviation has limiting upper average cost at least zero. The prescribed feedback makes the square identically zero and attains zero. No restriction to stationary deviations was used.

Finally, the norm of the first-moment vector is zero under $\pi_0$ and $r>0$ under $\pi_r$. Thus the two population-law distributions are distinct, even though averaging $\tau_\Phi m_r$ over phases gives the uniform one-particle unconditional distribution.
\end{proof}

\section{Why this lies outside the monotone class}
For $0<|\epsilon|<1$, let $m_\epsilon(dx)=(1+\epsilon\cos x)\ell(dx)$. Its cosine moment is $\epsilon/2$, its sine moment is zero, and direct integration gives
\[
 \int [F(x,m_\epsilon)-F(x,\ell)]\,[m_\epsilon-\ell](dx)
 =-\frac{\kappa\epsilon^2}{4}<0.
\]
The contribution of the squared sine term vanishes because $\int\sin^2x\cos x\,d\ell=0$. Thus the usual monotonicity inequality fails, not merely a strict version of it.

Additive common noise translates a nonuniform stationary profile without changing the magnitude of its first moment. The verification identity (2) shows that this persistence is compatible with optimal behavior, rather than being an arbitrary invariant solution of a population equation. The construction does not establish uniqueness within either branch, classify all stationary laws, or prove existence for arbitrary bounded nonmonotone couplings. It also does not refute positive results under monotonicity or different common-noise vector fields. Those are the remaining parts of the original problem.

\begin{thebibliography}{9}
\bibitem{CD} P. Cardaliaguet and F. Delarue (2022), ``Selected topics in mean field games,'' \emph{Proceedings of the International Congress of Mathematicians}. \href{https://doi.org/10.4171/ICM2022/17}{doi:10.4171/ICM2022/17}.
\bibitem{CMY} P. Cardaliaguet, R. Maillet and W. Yan (2025), ``On the long-time behavior of mean field game systems with a common noise,'' \href{https://arxiv.org/abs/2509.17443v1}{arXiv:2509.17443v1}. Theorems 3.9 and 6.6 concern the cited long-time and stationary theory; no general nonmonotone theorem is inferred from them.
\end{thebibliography}

\end{document}
